\subsection{Permanencia de las propiedades de los módulos por descenso fielmente plano.}

\begin{proposition}[2.5.1]
\label{IV.2.5.1-fr}
Sean $f:X\to Y$ un morfismo, $\mathscr F$ un $\mathcal O_X$-módulo
cuasicoherente, $g:Y'\to Y$ un morfismo fielmente plano,
$X'=X\times_Y Y'$, $f'=f_{(Y')}:X'\to Y'$,
$\mathscr F'=\mathscr F\otimes_{\mathcal O_Y}\mathcal O_{Y'}$. Para que
$\mathscr F$ sea plano \textup{(resp.} fielmente plano) con respecto a $f$, es necesario
y suficiente que $\mathscr F'$ sea plano \textup{(resp.} fielmente plano) con respecto
a $f'$.
\end{proposition}

Basta aplicar (2.2.10) reemplazando $X$, $Y'$, $Z$, $\mathscr F$ y
$\mathscr G'$ por $Y'$, $X$, $Y$, $\mathcal O_{Y'}$ y $\mathscr F$,
respectivamente; se concluye que, para que $\mathscr F$ sea plano (resp.
fielmente plano) con respecto a $f$, es necesario y suficiente que
$\mathscr F'$ sea plano (resp. fielmente plano) con respecto a $g\circ f'$.
Pero si $g':X'\to X$ es la proyección canónica, $g'$ es fielmente plano
(2.2.13) y se tiene $g\circ f'=f\circ g'$; por tanto, para que $\mathscr F'$
sea plano (resp. fielmente plano) con respecto a $f\circ g'$, es necesario y
suficiente que $\mathscr F$ lo sea con respecto a $f$ (2.2.11, (iii)).

\begin{proposition}[2.5.2]
\EGAReaderTarget{ega4.c04.prop.2-5-2}{IV.2.5.2-fr}
Sean $f:X'\to X$ un morfismo fielmente plano y cuasicompacto,
$\mathscr F$ un $\mathcal O_X$-módulo cuasicoherente y
$\mathscr F'=f^*(\mathscr F)$. Consideremos, para un módulo cuasicoherente,
la propiedad de ser:
\begin{enumerate}
\item[\textup{(i)}] de tipo finito;
\item[\textup{(ii)}] de presentación finita;
\item[\textup{(iii)}] localmente libre de tipo finito;
\item[\textup{(iv)}] localmente libre de rango $n$.
\end{enumerate}
Entonces, si $\boldsymbol{P}$ designa una de las propiedades precedentes, para que
$\mathscr F$ posea la propiedad $\boldsymbol{P}$, es necesario y suficiente que
$\mathscr F'$ la posea.
\end{proposition}

Para que un $\mathcal O_X$-módulo cuasicoherente sea localmente libre de tipo
finito, es necesario y suficiente que sea plano sobre $X$ y de presentación
finita (Bourbaki, \emph{Alg. comm.}, cap.~\textbf{II}, \S~5, núm.~2, cor.~2 del
teorema~1, habida cuenta de (2.1.2)); como $\mathscr F$ es plano sobre $X$ si
y solo si $\mathscr F'$ es plano sobre $X'$ en virtud de (2.5.1) (aplicado
tomando como $f$ la identidad), se ve que, para demostrar la proposición en el
caso (iii), basta haberla demostrado en los casos (i) y (ii); lo mismo ocurre
para demostrar (iv), pues
$f^*(\mathcal O_Y^n)=\mathcal O_X^n$, de modo que, si $\mathscr F$ y
$\mathscr F'$ son localmente libres de tipo finito y $x=f(x')$, el rango de
$\mathscr F$ en $x$ es igual al de $\mathscr F'$ en $x'$, y nuestra afirmación
se sigue de que $f$ es sobreyectivo. Para tratar los casos (i) y (ii),
\oldpage[IV-2]{23}
observemos que la cuestión es local sobre $X$, por lo que se puede suponer que
$X$ es afín; se sabe entonces por (2.2.12) que existe un esquema \emph{afín} $X''$ y
un morfismo fielmente plano $g:X''\to X'$ que es un isomorfismo local; por
consiguiente, es equivalente decir que $f^*(\mathscr F)$ posee la propiedad
$\boldsymbol{P}$, o que $g^*(f^*(\mathscr F))$ la posee. La cuestión se reduce así al
caso en que $X=\operatorname{Spec}(A)$ y $X'=\operatorname{Spec}(A')$; habida
cuenta de (2.2.3) y (\textbf{II}, 6.1.4.1), basta, pues, demostrar el

\begin{lemma}[2.5.3]
\label{IV.2.5.3-fr}
Sean $A$ un anillo, $A'$ una $A$-álgebra fielmente plana, $M$ un $A$-módulo y
$M'=M\otimes_A A'$. Para que $M$ sea de tipo finito \textup{(resp.} de presentación
finita), es necesario y suficiente que $M'$ lo sea.
\end{lemma}

Para la demostración, véase Bourbaki, \emph{Alg. comm.}, cap.~I, \S~3,
núm.~6, prop.~11.

